\documentclass[12pt]{article}

% ============================================================
% PACKAGES
% ============================================================

\usepackage{amsmath}      % Mathematical environments: align, equation, etc.
\usepackage{amsthm}       % Theorem, definition, proof environments
\usepackage{amsfonts}     % Mathematical fonts such as \mathbb
\usepackage{amssymb}      % Additional mathematical symbols
\usepackage{mathtools}    % Extensions to amsmath
\usepackage{geometry}     % Page margins
\usepackage{enumitem}     % Customize enumerate/itemize
\usepackage{mdframed}     % Framed boxes
\usepackage{tikz}         % Draw diagrams
\usepackage{hyperref}     % Clickable references and links

\geometry{
    letterpaper,
    margin=0.8in
}

\linespread{1.2}


% ============================================================
% NUMBER SYSTEMS
% ============================================================
% Instead of writing \mathbb{R} every time, use \RR.

\newcommand{\RR}{\mathbb{R}}   % Real numbers
\newcommand{\CC}{\mathbb{C}}   % Complex numbers
\newcommand{\NN}{\mathbb{N}}   % Natural numbers
\newcommand{\ZZ}{\mathbb{Z}}   % Integers
\newcommand{\QQ}{\mathbb{Q}}   % Rational numbers


% ============================================================
% COMMON MATHEMATICAL NOTATION
% ============================================================
% These commands are useful for frequently used notation.

\newcommand{\set}[1]{\left\{#1\right\}}
% Example:
%   \set{x \in \RR : x > 0}
% produces:
%   { x in R : x > 0 }

\newcommand{\abs}[1]{\left|#1\right|}
% Example:
%   \abs{x-y}
% produces:
%   |x-y|

\newcommand{\norm}[1]{\left\|#1\right\|}
% Example:
%   \norm{x}
% produces:
%   ||x||

\newcommand{\inner}[2]{\left\langle #1,#2\right\rangle}
% Example:
%   \inner{x}{y}
% produces:
%   <x,y>

\newcommand{\parens}[1]{\left(#1\right)}
% Example:
%   \parens{x+y}

\newcommand{\brac}[1]{\left[#1\right]}
% Example:
%   \brac{a,b}

\newcommand{\angles}[1]{\left\langle #1\right\rangle}
% Example:
%   \angles{x,y}


% ============================================================
% THEOREM / THEORY COMPONENTS
% ============================================================

% "definition" style is upright text.
% Use this for definitions, examples, and remarks.

\theoremstyle{definition}

% ------------------------------------------------------------
% DEFINITION
% ------------------------------------------------------------
% Use when introducing a formal mathematical definition.
%
% Usage:
%
% \begin{definition}[Metric]
% A metric on a set X is a function ...
% \end{definition}
%
% [Metric] is optional and gives the definition a name.

\newtheorem{definition}{Definition}[section]


% ------------------------------------------------------------
% EXAMPLE
% ------------------------------------------------------------
% Use to give a concrete example of a definition/theorem.
%
% Usage:
%
% \begin{example}
% Consider X = \RR with the usual distance.
% \end{example}

\newtheorem{example}[definition]{Example}


% ------------------------------------------------------------
% REMARK
% ------------------------------------------------------------
% Use for observations, comments, or small additional facts.
% It is NOT intended for major mathematical results.
%
% Usage:
%
% \begin{remark}
% This definition does not require X to be a subset of \RR^n.
% \end{remark}

\newtheorem{remark}[definition]{Remark}


% ============================================================
% THEOREM-LIKE COMPONENTS
% ============================================================

% "plain" style makes the mathematical statements italic.

\theoremstyle{plain}


% ------------------------------------------------------------
% THEOREM
% ------------------------------------------------------------
% Use for an important/general mathematical result.
%
% Usage:
%
% \begin{theorem}[Cauchy--Schwarz Inequality]
% For all x,y \in \RR^n,
% \[
% |\inner{x}{y}| \leq \norm{x}\norm{y}.
% \]
% \end{theorem}

\newtheorem{theorem}[definition]{Theorem}


% ------------------------------------------------------------
% PROPOSITION
% ------------------------------------------------------------
% Use for a useful mathematical result that is usually less
% central than a theorem.
%
% Usage:
%
% \begin{proposition}
% Every finite-dimensional subspace of \RR^n is closed.
% \end{proposition}

\newtheorem{proposition}[definition]{Proposition}


% ------------------------------------------------------------
% LEMMA
% ------------------------------------------------------------
% Use for an auxiliary result whose main purpose is to help
% prove another theorem/proposition.
%
% Usage:
%
% \begin{lemma}
% ...
% \end{lemma}

\newtheorem{lemma}[definition]{Lemma}




% ------------------------------------------------------------
% COROLLARY
% ------------------------------------------------------------
% Use for a result that follows relatively directly from a
% theorem/proposition that was just established.
%
% Usage:
%
% \begin{corollary}
% ...
% \end{corollary}

\newtheorem{corollary}[definition]{Corollary}


% ============================================================
% PROOF
% ============================================================
% The proof environment comes from amsthm.
%
% Usage:
%
% \begin{proof}
% Suppose that ...
%
% Therefore, ...
% \end{proof}
%
% LaTeX automatically adds the QED symbol at the end.


% ============================================================
% CUSTOM BOX: IDEA
% ============================================================
% Use this when you want to record intuition or the main idea
% behind a mathematical concept.
%
% Usage:
%
% \begin{idea}
% \textbf{Key idea.}
% A metric gives us a notion of distance...
% \end{idea}

\newmdenv[
    linewidth=1pt,
    skipabove=10pt,
    skipbelow=10pt
]{idea}


% ============================================================
% CUSTOM BOX: WARNING
% ============================================================
% Use this for common mistakes, distinctions, or things that
% are easy to confuse.
%
% Usage:
%
% \begin{warning}
% Do not confuse continuity with uniform continuity.
% \end{warning}

\newmdenv[
    linewidth=1pt,
    skipabove=10pt,
    skipbelow=10pt
]{warning}

% ============================================================
% CUSTOM BOX: NOTE
% ============================================================
% Use this for a short personal note, useful observation,
% connection between concepts, or explanation.
%
% Unlike "Remark", this is intended for informal notes.
%
% Usage:
%
% \begin{note}
% A sequence can be viewed as a function from $\NN$ to $\RR$.
% \end{note}

\newmdenv[
    linewidth=1pt,
    skipabove=10pt,
    skipbelow=10pt
]{note}


% ============================================================
% DOCUMENT
% ============================================================

\begin{document}


% ============================================================
% TITLE
% ============================================================

\begin{center}
    {\Large \textbf{Notations and The construction of $\mathbb{Q}$}}\\[0.5em]
    {\large Lecture 1}
\end{center}

% Automatically generate the table of contents.
\tableofcontents

\newpage


% ============================================================
% SECTION
% ============================================================
% Use \section{} for a major topic.
%
% Example:
%   \section{Metric Spaces}

\section{Notations}

% ============================================================
% SUBSECTION
% ============================================================
% Use \subsection{} when dividing a section into smaller topics.
%
% Example:
%   \subsection{Open Balls}

\begin{note}
Consider $x$, which is transcendental, if and only if $x$ is not the root of \textbf{any} non-zero polynomials, with integer coefficients.
\end{note}

\subsection{Set and Relations}

% ============================================================
% DEFINITION
% ============================================================

\begin{definition}[Set]
Set is a collection of (non-overlap) objects.
\end{definition}

\begin{note}
You can put anything in a set: numbers, even a smiley face, or the logo of "League of Legends".
\end{note}




% ============================================================
% REMARK
% ============================================================

\begin{remark}
We have 2 ways (at least) to describe a set, which is: list all of the objects in the set $B = \{a, b, c\}$ or tell the properties of objects in the set $A = \{x: x<2\}$.
\end{remark}

After that, the lecture introduced the notation of in ($\in$), not in ($\notin$), empty set ($\emptyset$), subset or equal ($\subseteq$), proper subset ($\subset$), and the condition for equal set ($A \subset B$ and $B \subset A$, then $A = B$).

\subsection{Operators on sets}

\begin{definition}[Union of sets]
$A \cup B$, is a set, which elements are $\in A$ or $\in B$.
\end{definition}

\begin{definition}[Intersection of sets]
$A \cap B$, is a set, which elements are $\in A$ and $\in B$.    
\end{definition}

\begin{definition}[Complement of sets]
$A^c$, is everything that not in A.
\end{definition}

\begin{definition}[Set minus]
$B \setminus A$, which is everything in B, but not in A.
\end{definition}

\begin{definition}[Set product]
$A \times B$, is a set of ordered pairs $(a, b)$, which $a \in A$ and $b \in B$.    
\end{definition}

\subsection{Relations}

\begin{definition}[Relation]
$R$, which is a relation, is a subset of $A \times B$. If $(a, b) \in R$ (we can say that: $a$ has relation $R$ with $b$) then we could write $aRb$.
\end{definition}

\begin{example}
$L$: loves, which is define on the set $P \times P$ (People $\times$ People). A pair $(Alex, Trix) \in L$, means that Alex is in love with Trix.
\end{example}

\begin{definition}[Equivalence relation]
An equivalence relation on a set $S$, is a relation $R$ over $S \times S$, such that:
\begin{itemize}
    \item \textbf{Reflexive} $aRa$ (means that $(a, a) \in R$).
    \item \textbf{Symmetric} $aRb \Rightarrow bRa$
    \item \textbf{Transitive} $aRb$ and $bRc$ $\Rightarrow$ $aRc$
\end{itemize}
\end{definition}

\begin{example}
$R$: 2 people in the same room is a equivalence relations.    
\end{example}

\begin{definition}[Function]
A function $F$ from $A$ to $B$ is a relation over $A \times B$ such that if $aFb$ and $aFb'$ then $b = b'$ (assign each $a$ to a unique $b$?).    
\end{definition}

\begin{remark}
We can just write $F(a)=b$ instead of $aFb$.
\end{remark}

\section{The Construction of Rational numbers}
\begin{note}
$\mathbb{Q}$, indeed, is a set, with the elements is "class" (to be more precise, "equivalence class"). But what is this "class"? These classes contains all pairs $(p,q) \in \mathbb{Z} \times (\mathbb{Z} \setminus \{0\})$  such that $(p,q)$ is in a equivalence relation with a given $(m,n)$ if $pn=qm$ and $q,n \neq 0$. We have: $[(m,n)]=\{(p,q) \in \mathbb{Z}\times (\mathbb{Z}\setminus {0}):(p,q) \sim (m,n)\}$. And note that, the important thing is, we want the arithmetic on $\mathbb{Z}$ is preserved on $\mathbb{Q}$.
\end{note}


\begin{definition}[Rational numbers $\mathbb{Q}$]
$\mathbb{Q} = \{\frac{p}{q}: p,q \in \mathbb{Z}, q \neq 0\}$ with $\frac{p}{q}$ is the equivalence class described above.
\end{definition}

\begin{note}
You may wonder if there exist both equivalence class of $(1,2)$ and $(3, 6)$. However, in set, the elements are non-overlap, and $[(1,2)] = [(3,6)]$. So, don't worry.
\end{note}

\begin{note}
About the equivalence relation for $(p,q)$ and $(m,n)$, it is a equivalence relation, satisfied all 3 conditions above.
\end{note}
\end{document}